\documentclass{article}
\usepackage[UTF8]{ctex}
\usepackage{amsmath,amssymb,amsfonts}
\usepackage{enumitem}
\usepackage{geometry}
\geometry{a4paper, margin=1in}
% === 错题本排版优化（不参与题目标识） ===
\setlength{\parindent}{0pt}
\setlength{\parskip}{0.35em plus 0.10em minus 0.08em}
\setlength{\emergencystretch}{3em}
\setlist{itemsep=0.25em, topsep=0.40em, parsep=0pt}
\allowdisplaybreaks[2]
% === 错题本排版优化结束 ===
\begin{document}
\section*{武汉理工大学2026年高等代数真题}
\section*{128．武汉理工大学2026年高等代数试卷}

\subsection{题目1}
1．计算 \(n\) 阶行列式

\[
D_{n}=\left|\begin{array}{ccccc}
1+b_{1} & 1 & 1 & \cdots & 1 \\
1 & 1+b_{2} & 1 & \cdots & 1 \\
1 & 1 & 1+b_{3} & \cdots & 1 \\
\vdots & \vdots & \vdots & & \vdots \\
1 & 1 & 1 & \cdots & 1+b_{n}
\end{array}\right|\left(b_{i} \neq 0, i=1,2, \cdots, n\right) .
\]

\subsection{题目2}
2．已知矩阵

\[
A=\left(\begin{array}{cccc}
2 & 4 & 0 & 0 \\
2 & 6 & 0 & 0 \\
0 & 0 & 0 & 4 \\
0 & 0 & -2 & 0
\end{array}\right) .
\]

若 \(\left[\left(\frac{1}{4} A\right)^{\bullet}\right]^{-1} B A^{-1}=\frac{1}{2} A B+6 E\) ，求 \(B\) 。

\subsection{题目3}
3．设向量

\[
\alpha_{1}=\left(\begin{array}{c}
a \\
a \\
a
\end{array}\right), \alpha_{2}=\left(\begin{array}{c}
b+1 \\
2 b+1 \\
b+1
\end{array}\right), \alpha_{2}=\left(\begin{array}{c}
2 \\
3 \\
b+4
\end{array}\right), \beta=\left(\begin{array}{c}
1 \\
1 \\
2 b+1
\end{array}\right) .
\]

\begin{itemize}
  \item[（1）] \(a, b\) 为何值时，\(\beta\) 不可以表示成 \(\alpha_{1}, \alpha_{2}, \alpha_{3}\) 的线性组合？
  \item[（2）] \(a, b\) 为何值时，\(\beta\) 可由 \(\alpha_{1}, \alpha_{2}, \alpha_{3}\) 线性表出？并写出线性表达式。
\end{itemize}

\subsection{题目4}
4．设 \(P\) 为数域，在 \(P^{2 \times 2}\) 中，令

\[
V_{1}=\left\{\left.\left(\begin{array}{cc}
x & -x \\
y & z
\end{array}\right) \right\rvert\, x, y, z \in P\right\}, V_{2}=\left\{\left.\left(\begin{array}{cc}
a & b \\
-a & c
\end{array}\right) \right\rvert\, a, b, c \in P\right\} .
\]

\begin{itemize}
  \item[（1）] 判断 \(V_{1}, V_{2}\) 是否为 \(P^{2 \times 2}\) 的子空间，并说明理由．
  \item[（2）] 分别求 \(V_{1}+V_{2}, V_{1} \cap V_{2}\) 的维数和一组基。
\end{itemize}

\subsection{题目5}
5．已知 \(\mathbb{R}^{3}\) 的线性变换 \(\sigma\) 对基

\[
\varepsilon_{1}=\left(\begin{array}{c}
-1 \\
0 \\
2
\end{array}\right), \varepsilon_{2}=\left(\begin{array}{l}
0 \\
1 \\
1
\end{array}\right), \varepsilon_{3}=\left(\begin{array}{c}
3 \\
-1 \\
-6
\end{array}\right)
\]

的像为

\[
\sigma\left(\varepsilon_{1}\right)=\left(\begin{array}{c}
-1 \\
0 \\
1
\end{array}\right), \sigma\left(\varepsilon_{2}\right)=\left(\begin{array}{c}
0 \\
-1 \\
-2
\end{array}\right), \sigma\left(\varepsilon_{3}\right)=\left(\begin{array}{c}
-1 \\
-1 \\
3
\end{array}\right) .
\]

\begin{itemize}
  \item[（1）] 求 \(\sigma\) 在基 \(\varepsilon_{1}, \varepsilon_{2}, \varepsilon_{3}\) 下的矩阵。
  \item[（2）] 设 \(X=\left(\begin{array}{l}1 \\ 1 \\ 1\end{array}\right)\) ，求 \(\sigma(X)\) ．
  \item[（3）] 已知 \(\sigma(Y)\) 在基 \(\varepsilon_{1}, \varepsilon_{2}, \varepsilon_{3}\) 下的坐标向量为 \(\left(\begin{array}{c}2 \\ -4 \\ -2\end{array}\right)\) ，求 \(Y\) ．
\end{itemize}

\subsection{题目6}
6．已知实二次型 \(f\left(x_{1}, x_{2}, x_{3}\right)=t\left(x_{1}^{2}+x_{2}^{2}+x_{3}^{2}\right)+4 x_{1} x_{2}+4 x_{1} x_{3}-4 x_{2} x_{3}\) ．

\begin{itemize}
  \item[（1）] 若 \(f\) 正定，求 \(t\) 的取值范围．
  \item[（2）] \(t=2\) 时，求 \(f\) 的规范形，并写出所作的非退化线性替换。
\end{itemize}

\subsection{题目7}
7．用 \(J\) 表示元素全为 1 的 \(n\) 阶方阵 \((n \geq 2)\) ，设 \(f(x)=a+b x \in \mathbb{Q}[x]\) ，令 \(A=f(J)\) 。

\begin{itemize}
  \item[（1）] 求 \(J\) 的全部特征值及特征向量。
  \item[（2）] 求 \(A\) 的所有特征子空间。
  \item[（3）] 问 \(A\) 是否可对角化？若可以，求可逆矩阵 \(P \in \mathbb{Q}^{n \times n}\) ，使得 \(P^{-1} A P\) 为对角阵．
\end{itemize}

\subsection{题目8}
8．设 \(\sigma, \tau\) 为 \(n\) 维线性空间 \(V\) 上的线性变换，满足 \(\sigma^{2}=\sigma\) 。证明：

\begin{itemize}
  \item[（1）] \(V=\operatorname{Im} \sigma \oplus \operatorname{Ker} \sigma\) ．
  \item[（2）] \(\operatorname{Im} \sigma, \operatorname{Ker} \sigma\) 均为 \(\tau\) 的不变子空间当且仅当 \(\sigma \tau=\tau \sigma\) ．
\end{itemize}

\subsection{题目9}
9．设 \(A, B\) 为 \(n\) 阶实对称正定矩阵，证明：\(|A+B|>|A|+|B|\) 。
\end{document}
